\documentclass[11pt,a4paper]{article}
\usepackage[utf8]{inputenc}
\usepackage[T1]{fontenc}
\usepackage{amsmath,amssymb,amsthm}
\usepackage{bm}
\usepackage{geometry}
\usepackage{booktabs}
\usepackage{algorithm}
\usepackage{algorithmic}
\usepackage{hyperref}

\geometry{margin=2.5cm}

\newcommand{\R}{\mathbb{R}}
\newcommand{\norm}[1]{\left\| #1 \right\|}
\newcommand{\wnorm}[2]{\left\| #1 \right\|_{#2}^2}
\newcommand{\argmin}{\operatorname*{arg\,min}}
\newcommand{\diag}{\operatorname{diag}}
\newcommand{\nullspace}{\operatorname{null}}

\title{OCS2 Legged Robot Whole-Body Control: \\Mathematical Problem Formulation}
\author{Note based on analysis of \texttt{legged\_control} WBC notes}
\date{\today}

\begin{document}
\maketitle
\tableofcontents
\newpage

%=============================================================================

\section{Overview}
%=============================================================================

The OCS2 legged robot controller is split into two optimization layers:
\begin{itemize}
	\item \textbf{Centroidal MPC}: a switched-system optimal control problem over a finite horizon. It plans the base motion, joint reference motion, contact forces, and contact mode sequence.
	\item \textbf{Whole-Body Control (WBC)}: a single-step inverse-dynamics quadratic program solved at the low-level control rate. It converts the MPC policy into dynamically consistent joint torques.
\end{itemize}

The reason for this separation is model fidelity and time scale. The MPC uses a reduced centroidal model to reason about future motion. The WBC uses the current full rigid-body model to produce torque commands that satisfy instantaneous dynamics, contact constraints, friction limits, and actuator limits.

\begin{center}
	\begin{tabular}{lll}
		\toprule
		& \textbf{MPC} & \textbf{WBC} \\
		\midrule
		Model & Centroidal / SRBD & Full floating-base rigid body model \\
		Rate & $\approx 50$ Hz & $\approx 400$ Hz \\
		Horizon & $1$ s receding horizon & Single control instant \\
		Decision variables & Contact forces, joint velocities & Base acceleration, contact forces, torques \\
		Output & Policy $\pi^*(t)$ & Motor torque command $\bm{\tau}_{cmd}$ \\
		\bottomrule
	\end{tabular}
\end{center}

At every WBC tick, the controller evaluates the latest MPC policy at the current time and current observation:
\begin{equation}
	(\bm{x}^*(t),\bm{u}^*(t),\sigma^*(t))
	=
	\operatorname{evaluatePolicy}\!\left(\pi^*,t,\hat{\bm{x}}(t)\right),
\end{equation}
where $\bm{x}^*$ is the optimized centroidal state, $\bm{u}^*$ is the optimized MPC input, and $\sigma^*$ is the planned contact mode. The WBC then solves a QP conditioned on this instantaneous reference.

%=============================================================================

\section{Interface Between MPC and WBC}
%=============================================================================

\subsection{MPC Policy Output}

The MPC policy evaluated at time $t$ provides:
\begin{equation}
	\bm{x}^*(t)
	=
	\begin{bmatrix}
		\dot{\bm{p}}_{\text{com}}^* \\
		\bm{\omega}_{\text{com}}^* \\
		\bm{r}_b^* \\
		\bm{\theta}_{zyx}^* \\
		\bm{q}_j^*
	\end{bmatrix},
	\qquad
	\bm{u}^*(t)
	=
	\begin{bmatrix}
		\bm{F}^{\text{mpc}} \\
		\dot{\bm{q}}_j^*
	\end{bmatrix}.
\end{equation}

Here $\bm{F}^{\text{mpc}} \in \R^{12}$ is the stacked ground reaction force reference:
\begin{equation}
	\bm{F}^{\text{mpc}}
	=
	\begin{bmatrix}
		\bm{f}_{\text{LF}}^{\text{mpc}} \\
		\bm{f}_{\text{RF}}^{\text{mpc}} \\
		\bm{f}_{\text{LH}}^{\text{mpc}} \\
		\bm{f}_{\text{RH}}^{\text{mpc}}
	\end{bmatrix},
	\qquad
	\bm{f}_i^{\text{mpc}} \in \R^3.
\end{equation}

The contact mode $\sigma^*(t)\in\{0,\ldots,15\}$ is converted into contact flags:
\begin{equation}
	c_i(t)
	=
	\begin{cases}
		1 & \text{if foot } i \text{ is in stance},\\
		0 & \text{if foot } i \text{ is in swing}.
	\end{cases}
\end{equation}

\subsection{Reference Quantities Used by WBC}

The WBC does not use the full MPC horizon directly. It extracts only the instantaneous reference values:
\begin{align}
	\bm{q}_j^{des}      &= \operatorname{jointAngles}(\bm{x}^*),\\
	\dot{\bm{q}}_j^{des}&= \operatorname{jointVelocities}(\bm{u}^*),\\
	\bm{p}_{f,i}^{des}  &= \operatorname{FK}_i(\bm{r}_b^*,\bm{\theta}_{zyx}^*,\bm{q}_j^*),\\
	\dot{\bm{p}}_{f,i}^{des}
	&= \bm{J}_i(\bm{q}^{des})\bm{v}^{des},\\
	\bm{F}^{des}        &= \bm{F}^{\text{mpc}}.
\end{align}

The WBC also builds a desired base acceleration from the desired centroidal momentum rate. This is the main mathematical bridge between the centroidal MPC and the full-body inverse dynamics QP.

%=============================================================================

\section{Rigid-Body Model Used by WBC}
%=============================================================================

\subsection{Generalized Coordinates and Velocities}

The WBC uses the floating-base rigid-body state in Pinocchio coordinates:
\begin{equation}
	\bm{q}
	=
	\begin{bmatrix}
		\bm{r}_b\\
		\bm{\theta}_{zyx}\\
		\bm{q}_j
	\end{bmatrix}
	\in \R^{18},
	\qquad
	\bm{v}
	=
	\begin{bmatrix}
		\bm{v}_b\\
		\bm{\omega}_b\\
		\dot{\bm{q}}_j
	\end{bmatrix}
	\in \R^{18}.
\end{equation}

At every control tick, the measured state is converted into $(\bm{q},\bm{v})$, and Pinocchio computes:
\begin{equation}
	(\bm{q},\bm{v})
	\longmapsto
	\left(
	\bm{M}(\bm{q}),\;
	\bm{n}(\bm{q},\bm{v}),\;
	\bm{J}_i(\bm{q}),\;
	\dot{\bm{J}}_i(\bm{q},\bm{v})
	\right),
\end{equation}
where $\bm{M}$ is the mass matrix, $\bm{n}$ is the nonlinear effects vector, and $\bm{J}_i$ is the translational Jacobian of foot $i$.

\subsection{Floating-Base Equations of Motion}

The full floating-base dynamics are:
\begin{equation}
	\bm{M}(\bm{q})\dot{\bm{v}}
	+
	\bm{n}(\bm{q},\bm{v})
	=
	\bm{S}^{\top}\bm{\tau}
	+
	\bm{J}_c(\bm{q})^{\top}\bm{F},
\end{equation}
where:
\begin{itemize}
	\item $\bm{\tau}\in\R^{12}$ is the actuated joint torque vector,
	\item $\bm{F}\in\R^{12}$ is the stacked contact force vector,
	\item $\bm{J}_c\in\R^{12\times 18}$ is the stacked foot Jacobian,
	\item $\bm{S} = [\bm{0}_{12\times 6}\;\bm{I}_{12}]$ selects actuated joints.
\end{itemize}

Because the floating base is unactuated, the first six rows are the most important hard consistency condition. Partition:
\begin{equation}
	\bm{M}
	=
	\begin{bmatrix}
		\bm{M}_b\\
		\bm{M}_j
	\end{bmatrix},
	\qquad
	\bm{n}
	=
	\begin{bmatrix}
		\bm{n}_b\\
		\bm{n}_j
	\end{bmatrix},
	\qquad
	\bm{J}_c
	=
	\begin{bmatrix}
		\bm{J}_{c,b} & \bm{J}_{c,j}
	\end{bmatrix},
\end{equation}
where $\bm{M}_b\in\R^{6\times18}$ and $\bm{n}_b\in\R^6$ are the floating-base rows. The floating-base equation is:
\begin{equation}
	\boxed{
		\bm{M}_b
		\begin{bmatrix}
			\dot{\bm{v}}_b\\
			\ddot{\bm{q}}_j
		\end{bmatrix}
		-
		\bm{J}_{c,b}^{\top}\bm{F}
		=
		-\bm{n}_b
	}
	\label{eq:floating_base_eom}
\end{equation}

In the WBC QP, $\dot{\bm{v}}_b$ is optimized directly. The joint acceleration $\ddot{\bm{q}}_j$ is obtained from the desired joint velocity update:
\begin{equation}
	\ddot{\bm{q}}_j^{des}
	\approx
	\frac{\dot{\bm{q}}_j^{des}(t)-\dot{\bm{q}}_j^{last}}{\Delta t}.
\end{equation}
Thus Eq.~\eqref{eq:floating_base_eom} becomes a linear equality in the QP variables $\dot{\bm{v}}_b$ and $\bm{F}$.

%=============================================================================

\section{WBC Decision Variables}
%=============================================================================

The WBC decision variable is:
\begin{equation}
	\boxed{
		\bm{z}
		=
		\begin{bmatrix}
			\dot{\bm{v}}_b\\
			\bm{F}\\
			\bm{\tau}
		\end{bmatrix}
		=
		\begin{bmatrix}
			\ddot{\bm{p}}_b\\
			\dot{\bm{\omega}}_b\\
			\bm{f}_{\text{LF}}\\
			\bm{f}_{\text{RF}}\\
			\bm{f}_{\text{LH}}\\
			\bm{f}_{\text{RH}}\\
			\bm{\tau}
		\end{bmatrix}
		\in \R^{30}
	}
\end{equation}

The optimized variables have distinct roles:
\begin{itemize}
	\item $\dot{\bm{v}}_b\in\R^6$: base linear and angular acceleration,
	\item $\bm{F}\in\R^{12}$: all candidate foot contact forces,
	\item $\bm{\tau}\in\R^{12}$: joint torques sent to the low-level command interface.
\end{itemize}

%=============================================================================

\section{Hard Constraints}
\label{sec:hard_constraints}
%=============================================================================

Hard constraints are constraints that the WBC must satisfy. They encode physics, contact consistency, actuator limits, and friction limits.

\subsection{Floating-Base Dynamics Constraint}

Let the full generalized acceleration be:
\begin{equation}
	\dot{\bm{v}}
	=
	\begin{bmatrix}
		\dot{\bm{v}}_b\\
		\ddot{\bm{q}}_j^{des}
	\end{bmatrix}.
\end{equation}
Using the floating-base rows of the rigid-body dynamics gives:
\begin{equation}
	\bm{A}_{eom}\bm{z} = \bm{b}_{eom},
\end{equation}
with
\begin{equation}
	\bm{A}_{eom}
	=
	\begin{bmatrix}
		\bm{M}_{bb} & -\bm{J}_{c,b}^{\top} & \bm{0}_{6\times12}
	\end{bmatrix},
	\qquad
	\bm{b}_{eom}
	=
	-\bm{n}_b - \bm{M}_{bj}\ddot{\bm{q}}_j^{des}.
\end{equation}

Here $\bm{M}_{bb}\in\R^{6\times6}$ and $\bm{M}_{bj}\in\R^{6\times12}$ are the floating-base columns of $\bm{M}_b$:
\begin{equation}
	\bm{M}_b
	=
	\begin{bmatrix}
		\bm{M}_{bb} & \bm{M}_{bj}
	\end{bmatrix}.
\end{equation}

\subsection{Torque Limit Constraint}

Each actuator torque is bounded:
\begin{equation}
	\boxed{
		-\bm{\tau}_{max}
		\leq
		\bm{\tau}
		\leq
		\bm{\tau}_{max}
	}
\end{equation}
This is represented as a box inequality on the torque block of $\bm{z}$.

\subsection{Friction Pyramid Constraint}

For a stance foot $i$ with $c_i=1$, the contact force
$\bm{f}_i=[f_{x,i},f_{y,i},f_{z,i}]^\top$ must lie inside a linearized Coulomb cone:
\begin{equation}
	\boxed{
		\begin{bmatrix}
			0 & 0 & -1\\
			1 & 0 & -\mu\\
			-1 & 0 & -\mu\\
			0 & 1 & -\mu\\
			0 & -1 & -\mu
		\end{bmatrix}
		\bm{f}_i
		\leq
		\bm{0}
	}
	\label{eq:friction_pyramid}
\end{equation}
Equivalently:
\begin{equation}
	f_{z,i}\geq 0,\qquad
	|f_{x,i}|\leq \mu f_{z,i},\qquad
	|f_{y,i}|\leq \mu f_{z,i}.
\end{equation}

For a swing foot $i$ with $c_i=0$, the contact force must be zero:
\begin{equation}
	\boxed{
		\bm{f}_i = \bm{0}
	}
\end{equation}
This prevents the QP from using nonphysical ground reaction forces on non-contacting feet.

\subsection{No-Contact-Motion Constraint}

For a stance foot, the foot acceleration should be zero in the world frame:
\begin{equation}
	\ddot{\bm{p}}_{f,i}
	=
	\bm{J}_i\dot{\bm{v}}
	+
	\dot{\bm{J}}_i\bm{v}
	=
	\bm{0}.
\end{equation}
With the optimized base acceleration and fixed desired joint acceleration:
\begin{equation}
	\bm{J}_{i,b}\dot{\bm{v}}_b
	=
	-\dot{\bm{J}}_i\bm{v}
	-
	\bm{J}_{i,j}\ddot{\bm{q}}_j^{des}.
\end{equation}
Stacking all stance feet gives:
\begin{equation}
	\bm{A}_{stance}\bm{z}=\bm{b}_{stance}.
\end{equation}

%=============================================================================

\section{Soft Tasks}
\label{sec:soft_tasks}
%=============================================================================

Soft tasks define what the WBC should track when the hard constraints leave degrees of freedom. They are written as least-squares objectives:
\begin{equation}
	\bm{A}_{task}\bm{z} \approx \bm{b}_{task}.
\end{equation}

\subsection{Base Acceleration Tracking}

The desired base acceleration comes from the MPC centroidal momentum evolution. Let the centroidal momentum be:
\begin{equation}
	\bm{h}
	=
	\bm{A}_g(\bm{q})\bm{v},
	\qquad
	\dot{\bm{h}}
	=
	\bm{A}_g(\bm{q})\dot{\bm{v}}
	+
	\dot{\bm{A}}_g(\bm{q},\bm{v})\bm{v}.
\end{equation}
Partition the centroidal momentum matrix:
\begin{equation}
	\bm{A}_g
	=
	\begin{bmatrix}
		\bm{A}_{g,b} & \bm{A}_{g,j}
	\end{bmatrix},
\end{equation}
where $\bm{A}_{g,b}\in\R^{6\times6}$ maps base velocity to centroidal momentum and $\bm{A}_{g,j}\in\R^{6\times12}$ maps joint velocity.

The MPC provides normalized centroidal momentum. If
$\dot{\bm{h}}_{norm}^{des}$ is the desired normalized centroidal momentum rate, then:
\begin{equation}
	\dot{\bm{h}}^{des}
	=
	m\dot{\bm{h}}_{norm}^{des}.
\end{equation}
Solving the centroidal relation for desired base acceleration gives:
\begin{equation}
	\boxed{
		\dot{\bm{v}}_b^{des}
		=
		\bm{A}_{g,b}^{-1}
		\left(
			\dot{\bm{h}}^{des}
			-
			\dot{\bm{A}}_g\bm{v}
			-
			\bm{A}_{g,j}\ddot{\bm{q}}_j^{des}
		\right)
	}
	\label{eq:base_acc_from_centroidal}
\end{equation}

The corresponding least-squares task is:
\begin{equation}
	\bm{A}_{base}\bm{z} \approx \bm{b}_{base},
	\qquad
	\bm{A}_{base}
	=
	\begin{bmatrix}
		\bm{I}_6 & \bm{0}_{6\times12} & \bm{0}_{6\times12}
	\end{bmatrix},
	\qquad
	\bm{b}_{base}
	=
	\dot{\bm{v}}_b^{des}.
\end{equation}

\subsection{Swing Leg Tracking}

For each swing foot, the WBC uses task-space PD to generate a desired foot acceleration:
\begin{equation}
	\bm{a}_{f,i}^{des}
	=
	\bm{K}_{p,f}
	\left(
		\bm{p}_{f,i}^{des}-\bm{p}_{f,i}
	\right)
	+
	\bm{K}_{d,f}
	\left(
		\dot{\bm{p}}_{f,i}^{des}-\dot{\bm{p}}_{f,i}
	\right).
\end{equation}

The actual foot acceleration is:
\begin{equation}
	\ddot{\bm{p}}_{f,i}
	=
	\bm{J}_{i,b}\dot{\bm{v}}_b
	+
	\bm{J}_{i,j}\ddot{\bm{q}}_j^{des}
	+
	\dot{\bm{J}}_i\bm{v}.
\end{equation}
Therefore the swing tracking task is:
\begin{equation}
	\bm{J}_{i,b}\dot{\bm{v}}_b
	\approx
	\bm{a}_{f,i}^{des}
	-
	\dot{\bm{J}}_i\bm{v}
	-
	\bm{J}_{i,j}\ddot{\bm{q}}_j^{des}.
\end{equation}
In QP matrix form:
\begin{equation}
	\bm{A}_{swing,i}\bm{z}
	\approx
	\bm{b}_{swing,i},
	\qquad
	\bm{A}_{swing,i}
	=
	\begin{bmatrix}
		\bm{J}_{i,b} & \bm{0}_{3\times12} & \bm{0}_{3\times12}
	\end{bmatrix}.
\end{equation}

\subsection{Contact Force Tracking}

The MPC computes a feedforward ground reaction force distribution. The WBC tracks this force while still respecting full-body dynamics and friction:
\begin{equation}
	\bm{F}
	\approx
	\bm{F}^{des}.
\end{equation}
Only stance forces are meaningful, so the task is contact-gated:
\begin{equation}
	\bm{C}_F\bm{F}
	\approx
	\bm{C}_F\bm{F}^{des},
	\qquad
	\bm{C}_F
	=
	\diag(c_{\text{LF}}\bm{I}_3,\;
	      c_{\text{RF}}\bm{I}_3,\;
	      c_{\text{LH}}\bm{I}_3,\;
	      c_{\text{RH}}\bm{I}_3).
\end{equation}
The QP task is:
\begin{equation}
	\bm{A}_{force}\bm{z}
	\approx
	\bm{b}_{force},
	\qquad
	\bm{A}_{force}
	=
	\begin{bmatrix}
		\bm{0}_{12\times6} & \bm{C}_F & \bm{0}_{12\times12}
	\end{bmatrix},
	\qquad
	\bm{b}_{force}
	=
	\bm{C}_F\bm{F}^{des}.
\end{equation}

%=============================================================================

\section{Weighted WBC Formulation}
\label{sec:weighted_wbc}
%=============================================================================

The default WBC implementation solves one constrained QP. Hard constraints are enforced exactly, while soft tasks are combined into a weighted least-squares cost.

\subsection{Weighted Task Stack}

Define the weighted task stack:
\begin{equation}
	\bm{A}_w
	=
	\begin{bmatrix}
		\bm{W}_{base}\bm{A}_{base}\\
		\bm{W}_{swing}\bm{A}_{swing}\\
		\bm{W}_{force}\bm{A}_{force}
	\end{bmatrix},
	\qquad
	\bm{b}_w
	=
	\begin{bmatrix}
		\bm{W}_{base}\bm{b}_{base}\\
		\bm{W}_{swing}\bm{b}_{swing}\\
		\bm{W}_{force}\bm{b}_{force}
	\end{bmatrix}.
\end{equation}

The soft objective is:
\begin{equation}
	\min_{\bm{z}}
	\frac{1}{2}
	\norm{\bm{A}_w\bm{z}-\bm{b}_w}^2.
\end{equation}
Expanding into standard QP form gives:
\begin{equation}
	\min_{\bm{z}}
	\frac{1}{2}\bm{z}^{\top}\bm{H}\bm{z}
	+
	\bm{g}^{\top}\bm{z},
	\qquad
	\bm{H}=\bm{A}_w^{\top}\bm{A}_w,
	\qquad
	\bm{g}=-\bm{A}_w^{\top}\bm{b}_w.
\end{equation}

\subsection{Constraint Stack}

All equality and inequality constraints are stacked into:
\begin{equation}
	\bm{A}_{eq}\bm{z}=\bm{b}_{eq},
	\qquad
	\bm{D}\bm{z}\leq \bm{f}.
\end{equation}
For qpOASES, they are represented as lower and upper bounds on a single constraint matrix:
\begin{equation}
	\bm{A}_{qp}
	=
	\begin{bmatrix}
		\bm{A}_{eq}\\
		\bm{D}
	\end{bmatrix},
	\qquad
	\bm{\ell}_{A}
	=
	\begin{bmatrix}
		\bm{b}_{eq}\\
		-\infty
	\end{bmatrix},
	\qquad
	\bm{u}_{A}
	=
	\begin{bmatrix}
		\bm{b}_{eq}\\
		\bm{f}
	\end{bmatrix}.
\end{equation}

\subsection{QP Problem}

The weighted WBC problem is:
\begin{equation}
	\boxed{
		\begin{aligned}
			\bm{z}^*
			=
			\argmin_{\bm{z}}\quad
			&
			\frac{1}{2}\bm{z}^{\top}\bm{H}\bm{z}
			+
			\bm{g}^{\top}\bm{z}\\
			\text{s.t.}\quad
			&
			\bm{A}_{eq}\bm{z}=\bm{b}_{eq},\\
			&
			\bm{D}\bm{z}\leq\bm{f},\\
			&
			-\bm{\tau}_{max}\leq\bm{\tau}\leq\bm{\tau}_{max}.
		\end{aligned}
	}
\end{equation}

This formulation is computationally efficient because only one QP is solved per WBC tick. The trade-off is that the relative importance of base tracking, swing tracking, and force tracking depends on weight tuning.

%=============================================================================

\section{Hierarchical WBC and HoQP}
\label{sec:hierarchical_wbc}
%=============================================================================

The codebase also contains a hierarchical WBC implementation. Instead of combining all soft tasks with weights, it solves a sequence of QPs in priority order.

\subsection{Task Priority Structure}

The tasks are grouped as:
\begin{align}
	\mathcal{T}_0
	&=
	\mathcal{T}_{eom}
	+
	\mathcal{T}_{torque}
	+
	\mathcal{T}_{friction}
	+
	\mathcal{T}_{stance},\\
	\mathcal{T}_1
	&=
	\mathcal{T}_{base}
	+
	\mathcal{T}_{swing},\\
	\mathcal{T}_2
	&=
	\mathcal{T}_{force}.
\end{align}
The priority order is:
\begin{equation}
	\mathcal{T}_0
	\succ
	\mathcal{T}_1
	\succ
	\mathcal{T}_2.
\end{equation}

Thus physical feasibility and contact consistency have highest priority, base and swing motion tracking have second priority, and MPC contact force tracking has the lowest priority.

\subsection{Null-Space Recursion}

Assume the solution after the previous higher-priority levels is $\bm{z}_{prev}$, and the null-space basis that preserves all previous equality tasks is $\bm{Z}_{prev}$. A lower-priority update is restricted to:
\begin{equation}
	\bm{z}
	=
	\bm{z}_{prev}
	+
	\bm{Z}_{prev}\Delta\bm{y}.
\end{equation}

For a current equality task
\begin{equation}
	\bm{A}\bm{z}=\bm{b},
\end{equation}
the HoQP level minimizes:
\begin{equation}
	\min_{\Delta\bm{y}}
	\frac{1}{2}
	\bigl\|
		\bm{A}
		\bigl(
			\bm{z}_{prev}+\bm{Z}_{prev}\Delta\bm{y}
		\bigr)
		-
		\bm{b}
	\bigr\|^2.
\end{equation}
The QP Hessian and linear term are:
\begin{equation}
	\bm{H}
	=
	\bm{Z}_{prev}^{\top}\bm{A}^{\top}\bm{A}\bm{Z}_{prev},
	\qquad
	\bm{c}
	=
	(\bm{A}\bm{Z}_{prev})^{\top}
	(\bm{A}\bm{z}_{prev}-\bm{b}).
\end{equation}

After solving, the solution is updated:
\begin{equation}
	\bm{z}_{new}
	=
	\bm{z}_{prev}
	+
	\bm{Z}_{prev}\Delta\bm{y}^*.
\end{equation}
The null-space basis is also updated:
\begin{equation}
	\boxed{
		\bm{Z}_{new}
		=
		\bm{Z}_{prev}
		\nullspace(\bm{A}\bm{Z}_{prev})
	}
\end{equation}

\subsection{Slack Variables for Inequalities}

When a level includes inequalities,
\begin{equation}
	\bm{D}\bm{z}\leq\bm{f},
\end{equation}
the HoQP formulation introduces slack variables $\bm{s}$ so that the level remains feasible:
\begin{equation}
	\bm{D}
	\left(
		\bm{z}_{prev}+\bm{Z}_{prev}\Delta\bm{y}
	\right)
	-
	\bm{s}
	\leq
	\bm{f}.
\end{equation}
A typical level problem is:
\begin{equation}
	\boxed{
		\begin{aligned}
			\min_{\Delta\bm{y},\bm{s}}\quad
			&
			\frac{1}{2}
			\norm{
				\bm{A}
				(\bm{z}_{prev}+\bm{Z}_{prev}\Delta\bm{y})
				-
				\bm{b}
			}^2
			+
			\frac{1}{2}\norm{\bm{s}}^2\\
			\text{s.t.}\quad
			&
			\bm{D}
			(\bm{z}_{prev}+\bm{Z}_{prev}\Delta\bm{y})
			-
			\bm{s}
			\leq
			\bm{f}.
		\end{aligned}
	}
\end{equation}

The final hierarchical solution can be written compactly as:
\begin{equation}
	\bm{z}^{*}
	=
	\operatorname{HoQP}
	\left(
		\mathcal{T}_2,\;
		\operatorname{HoQP}
		\left(
			\mathcal{T}_1,\;
			\operatorname{HoQP}(\mathcal{T}_0)
		\right)
	\right).
\end{equation}

%=============================================================================

\section{Command Extraction}
%=============================================================================

After solving the WBC QP, the torque command is extracted from the last block of the optimizer solution:
\begin{equation}
	\bm{z}^{*}
	=
	\begin{bmatrix}
		\dot{\bm{v}}_b^{*}\\
		\bm{F}^{*}\\
		\bm{\tau}^{*}
	\end{bmatrix},
	\qquad
	\boxed{
		\bm{\tau}_{wbc}
		=
		\bm{\tau}^{*}
	}
\end{equation}

The command sent to the hardware also includes joint position and velocity references from the MPC policy:
\begin{align}
	\bm{q}_{cmd} &= \bm{q}_j^{des},\\
	\dot{\bm{q}}_{cmd} &= \dot{\bm{q}}_j^{des},\\
	\bm{\tau}_{cmd} &= \bm{\tau}_{wbc}.
\end{align}

The low-level actuator interface may additionally apply fixed joint-space PD gains:
\begin{equation}
	\bm{\tau}_{motor}
	=
	\bm{\tau}_{cmd}
	+
	\bm{K}_{p,j}
	(\bm{q}_{cmd}-\bm{q}_j)
	+
	\bm{K}_{d,j}
	(\dot{\bm{q}}_{cmd}-\dot{\bm{q}}_j).
\end{equation}

Thus the PD terms are not the main whole-body controller. They are low-level stabilizing terms around the QP torque command, while the physically consistent torque distribution is computed by inverse-dynamics optimization.

%=============================================================================

\section{Complete Control Flow}
%=============================================================================

At each low-level control tick, the WBC pipeline is:
\begin{algorithm}
	\caption{MPC--WBC Low-Level Control Tick}
	\begin{algorithmic}[1]
		\STATE Receive current observation $(t,\hat{\bm{x}},\hat{\sigma})$
		\STATE Swap in the latest MPC policy if a new one is available
		\STATE Evaluate policy:
		$(\bm{x}^{*},\bm{u}^{*},\sigma^{*})
		\leftarrow
		\operatorname{evaluatePolicy}(\pi^*,t,\hat{\bm{x}})$
		\STATE Convert $\sigma^{*}$ to contact flags $c_i$
		\STATE Update measured rigid-body quantities:
		$\bm{M},\bm{n},\bm{J}_i,\dot{\bm{J}}_i$
		\STATE Update desired quantities:
		$\bm{q}_j^{des}$, $\dot{\bm{q}}_j^{des}$,
		$\bm{p}_{f,i}^{des}$, $\dot{\bm{p}}_{f,i}^{des}$,
		$\dot{\bm{v}}_b^{des}$, $\bm{F}^{des}$
		\STATE Build hard constraints:
		EoM, torque limits, friction pyramids, stance foot acceleration
		\STATE Build soft tasks:
		base acceleration, swing leg acceleration, contact force tracking
		\STATE Solve Weighted WBC QP or Hierarchical HoQP
		\STATE Extract $\bm{\tau}_{wbc}$ from $\bm{z}^{*}$
		\STATE Send $(\bm{q}_{cmd},\dot{\bm{q}}_{cmd},\bm{\tau}_{wbc})$ to hardware
	\end{algorithmic}
\end{algorithm}

\subsection{Conceptual Summary}

The MPC chooses \emph{what should happen} over the next horizon: body motion, contact schedule, joint motion reference, and approximate contact forces. The WBC chooses \emph{how to realize it now}: it computes torques that satisfy the current full-body dynamics and contact constraints while tracking the MPC references as closely as possible.

In mathematical form, the WBC is the instantaneous projection:
\begin{equation}
	(\bm{x}^{*},\bm{u}^{*},\sigma^{*})
	\quad
	\longmapsto
	\quad
	\bm{z}^{*}
	=
	\begin{bmatrix}
		\dot{\bm{v}}_b^{*}\\
		\bm{F}^{*}\\
		\bm{\tau}^{*}
	\end{bmatrix}
	\quad
	\longmapsto
	\quad
	\bm{\tau}_{motor}.
\end{equation}

\end{document}
